

\section{Background}
~{\vspace{-3ex}}
\begin{wrapfigure}{r}{0.55\textwidth}\vspace{-7.6ex}
  \centering
%   \begin{minipage}{0.6\textwidth}
% \centering
%   \includegraphics[width=0.95\textwidth]{figures/pipeline_ml.pdf}\vspace{-3ex}
%   \caption{{\small End-to-end robot learning.}}
  
%   \label{fig:end_to_end_pipeline}
%   \end{minipage}
%   \hspace{-3ex}\begin{minipage}{0.6\textwidth}
%     \centering
    % \vspace{1ex}
\resizebox{0.98\linewidth}{!}{
\begin{tikzpicture}[
  timeline/.style={ultra thick, draw=black!60, rounded corners=3pt},
    label/.style={align=center,  text width=2.3cm, inner sep=1pt},
    label2/.style={align=center,  text width=3cm, inner sep=1pt},
    event2022/.style={rounded corners=3pt, draw=black!70, fill=white, thick, minimum size=6mm, align=center, text=black},
    event2023/.style={rounded corners=3pt, draw=black!70, fill=gray!10, thick, minimum size=6mm,  align=center, text=black},
    event2024/.style={rounded corners=3pt, draw=black!70, fill=gray!50, thick, minimum size=6mm, align=center, text=black},
    event2025/.style={rounded corners=3pt, draw=black!70, fill=gray, thick, minimum size=6mm, align=center, text=black}
    ]

    % Main timeline path with reduced vertical spacing
    \draw[timeline,line width=6pt] (0,0) -- (15,0);
    \draw[-, line width=0.5pt, draw=white] (0,0) -- (15,0);
    % Arrow head at (0,0)
    \draw[-stealth, line width=6pt, draw=black!60] (0,0) -- (1,0);
    \draw[timeline,line width=6pt] (15,0) to[out=-0, in=0] (15,-1);
    \draw[-, line width=0.5pt, draw=white] (15,0) to[out=-0, in=0] (15,-1);
    \draw[timeline,line width=6pt] (15,-1) -- (0.5,-1);
    \draw[-, line width=0.5pt, draw=white] (15,-1) -- (0.5,-1);
    \draw[timeline,line width=6pt] (0.5,-1) to[out=-180, in=180] (0.5,-2.5);
    \draw[-, line width=0.5pt, draw=white] (0.5,-1) to[out=-180, in=180] (0.5,-2.5);
    \draw[timeline,line width=6pt] (0.5,-2.5) -- (15,-2.5);
    \draw[-, line width=0.5pt, draw=white] (0.5,-2.5) -- (15,-2.5);
    \draw[timeline,line width=6pt] (15,-2.5) to[out=-0, in=0] (15,-4);
    \draw[-, line width=0.5pt, draw=white] (15,-3) to[out=-0, in=0] (15,-4);
    \draw[timeline,line width=6pt] (15,-4) -- (0.5,-4); 
    \draw[-, line width=0.5pt, draw=white] (15,-4) -- (0.5,-4);
    \draw[timeline,line width=6pt] (0.5,-4) to[out=-180, in=180] (0.5,-5.5);
    \draw[-, line width=0.5pt, draw=white] (0.5,-4) to[out=-180, in=180] (0.5,-5.5);
    \draw[timeline,line width=6pt] (0.5,-5.5) -- (12.9,-5.5);
    \draw[-, line width=0.5pt, draw=white] (0.5,-5.5) -- (12.9,-5.5);
    \draw[-stealth, line width=6pt, draw=black!60] (12.9,-5.5) -- (14.1,-5.5);

    %%% Foundation Models %%%
    \node[event2022] (cliport) at (2.5,0) {\cite{shridhar2022cliport}};
    \node[label, above=3pt of cliport] {CLIPort};
    
    % \node[event2022] (gato) at (2.5,0) {\cite{reed2022generalist}};
    % \node[label, above=3pt of gato] {Gato};
    
    \node[event2022] (rt1) at (4,0) {\cite{rt-1}};
    \node[label, above=3pt of rt1] {RT-1};
    
    \node[event2023] (vima) at (5.5,0) {\cite{jiang2023vima}};
    \node[label, above=3pt of vima] {VIMA};
    
    \node[event2023] (act) at (7,0) {\cite{zhao2023learningFB}};
    \node[label, above=3pt of act] {ACT };
    
    \node[event2023] (rt2) at (8.5,0) {\cite{rt-2}};
    \node[label, above=3pt of rt2] {RT-2 };
    
    \node[event2023] (voxtposer) at (10,0) {\cite{huang2023voxposer}};
    \node[label, above=3pt of voxtposer] {VoxPoser };
    
    % \node[event2023] (diffpolicy) at (11.5,0) {\cite{chi2023diffusionpolicy}};
    % \node[label, above=3pt of diffpolicy] {Diffusion Policy };
    
    \node[event2024] (octo) at (12,0) {\cite{team2024octo}};
    \node[label, above=3pt of octo] {Octo };
    
    \node[event2024] (openvla) at (14,0) {\cite{kim2024openvla}};
    \node[label, above=3pt of openvla] {OpenVLA };

    %%% Specialized Applications %%%
    \node[event2024] (deer) at (1.1,-1) {\cite{yue2024deer}};
    \node[label, below=3pt of deer] {Deer-VLA };
    
    \node[event2024] (uni) at (2.9,-1) {\cite{zhang2024uni}};
    \node[label, below=3pt of uni] {Uni-NaVid };
    
    \node[event2024] (revla) at (4.6,-1) {\cite{dey2024revla}};
    \node[label, below=3pt of revla] {ReVLA };

    \node[event2024] (occllama) at (6.1,-1) {\cite{wei2024occllama}};
    \node[label, below=3pt of occllama] {Occllama };

    \node[event2024] (pi0) at (7.4,-1) {\cite{black2024pi_0}};
    \node[label, below=3pt of pi0] {Pi-01 };

    \node[event2024] (rdt) at (8.7,-1) {\cite{liu2024rdt}};
    \node[label, below=3pt of rdt] {RDT-1B };

    \node[event2024] (cogact) at (10.2,-1) {\cite{li2024cogact}};
    \node[label, below=3pt of cogact] {CogAct };

    \node[event2024] (edgevla) at (11.9,-1) {\cite{budzianowski20edgevla}};
    \node[label, below=3pt of edgevla] {EdgeVLA };

    \node[event2024] (showui) at (13.6,-1) {\cite{lin2024showui}};
    \node[label, below=3pt of showui] {ShowUI };

    %%% Robotics Focus %%%
    \node[event2024] (navila) at (1, -2.5) {\cite{cheng2024navila}};
    \node[label, below=3pt of navila] {NaviLa};

    \node[event2024] (quar) at (2.5, -2.5) {\cite{ding2024quar}};
    \node[label, below=3pt of quar] {Quar-VLA};

    \node[event2024] (bivla) at (4.4, -2.5) {\cite{gbagbe2024bi}};
    \node[label, below=3pt of bivla] {Bi-VLA};

    \node[event2024] (robomamba) at (6.3, -2.5) {\cite{liu2024robomamba}};
    \node[label, below=3pt of robomamba] {RoboMamba};

    \node[event2025] (otter) at (7.8, -2.5) {\cite{huang2025otter}};
    \node[label, below=3pt of otter] {Otter};

    \node[event2025] (pointvla) at (9.2, -2.5) {\cite{li2025pointvla}};
    \node[label, below=3pt of pointvla] {PointVLA};

    \node[event2025] (combatvla) at (11.3, -2.5) {\cite{chen2025combatvla}};
    \node[label, below=3pt of combatvla] {CombatVLA};

    \node[event2025] (hybridvla) at (13.5, -2.5) {\cite{liu2025hybridvla}};
    \node[label, below=3pt of hybridvla] {HybridVLA};

    %%% General Capabilities %%%
    \node[event2025] (covla) at (2.7, -4) {\cite{arai2025covla}}; % Changed from -6 to -5.2
    \node[label, below=3pt of covla] {CoVLA};
    % \node[event2025] (opendrivevla) at (2.7, -4) {\cite{zhou2025opendrivevla}}; % Changed from -6 to -5.2
    % \node[label, below=3pt of opendrivevla] {OpenDriveVLA};

    \node[event2025] (orion) at (4.6, -4) {\cite{fu2025orion}}; % Changed from -6 to -5.2
    \node[label, below=3pt of orion] {ORION};
    
    \node[event2025] (objectvla) at (6.3, -4) {\cite{zhu2025objectvla}}; % Changed from -6 to -5.2
    \node[label, below=3pt of objectvla] {ObjectVLA };

    \node[event2025] (conrft) at (8, -4) {\cite{chen2025conrft}}; % Changed from -6 to -5.2
    \node[label, below=3pt of conrft] {ConRFT};

    \node[event2025] (hirobot) at (9.5, -4) {\cite{shi2025hi}}; % Changed from -6 to -5.2
    \node[label, below=3pt of hirobot] {Hi Robot};

    \node[event2025] (tla) at (11, -4) {\cite{hao2025tla}}; % Changed from -6 to -5.2
    \node[label, below=3pt of tla] {TLA};
    
    \node[event2025] (racevla) at (12.5, -4) {\cite{serpiva2025racevla}}; % Changed from -6 to -5.2
    \node[label, below=3pt of racevla] {RaceVLA};

    \node[event2025] (dexvla) at (14.5, -4) {\cite{wen2025dexvla}}; % Changed from -6 to -5.2
    \node[label, below=3pt of dexvla] {DexVLA};
    %%% Newest Generalist Models (Q1--Q2 2025) %%%
    \node[event2024] (gr2) at (1.3, -4) {\cite{cheang2024gr2}};
    \node[label, below=3pt of gr2] {GR-2};

    \node[event2025] (spatialvla) at (1.3,-5.5) {\cite{qu2025spatialvla}};
    \node[label, below=3pt of spatialvla] {SpatialVLA};

    \node[event2025] (openvlaoft) at (3.7,-5.5) {\cite{openvla_oft}};
    \node[label2, below=3pt of openvlaoft] {OpenVLA-OFT};

    \node[event2025] (grootn1) at (6,-5.5) {\cite{gr00tn1}};
    \node[label, below=3pt of grootn1] {GR00T N1};

    \node[event2025] (geminirobotics) at (8.6,-5.5) {\cite{geminirobotics2025}};
    \node[label2, below=3pt of geminirobotics] {Gemini Robotics};

    \node[event2025] (pi05) at (10.8,-5.5) {\cite{pi05_2025}};
    \node[label, below=3pt of pi05] {Pi-0.5};

    \node[event2025] (worldvla) at (12.7,-5.5) {\cite{cen2025worldvla}};
    \node[label, below=3pt of worldvla] {WorldVLA};\end{tikzpicture}
}
\vspace{-1ex}
\caption{\scriptsize Recent robotic VLA models from {\fcolorbox{black}{white}{\scriptsize \text{2022}}} to {\fcolorbox{black}{gray}{\scriptsize \text{2025}}}.}\vspace{-1ex}
\label{fig:vla-timeline}
%   \end{minipage}
  \vspace{-2ex}
\end{wrapfigure}
% {\vspace{-1ex}} 
\BfPara{End-to-end Robotic Learning}
Building generalist
robots advances in multiple dimensions: 
\cib{1} the development and integration of vision, language, large action models, and multimodal fusion techniques into robotic learning, 
\cib{2} {\xo}-embodiment sensing and training using large, diverse datasets from across embodiments, enhanced transfer learning methods, and scalable data collection frameworks, and
\cib{3} the incorporation of real-time feedback, reinforcement learning on modular robotic models for various tasks.
High-capacity models have demonstrated an exceptional ability to process and learn diverse and extensive datasets, which is critical for end-to-end robotic learning. 
These models leverage their vast parameter space, \eg GPT-3, Jurassic-1, WuDao-2 with 175B+, Microsoft and Nvidia's Megatron-Turing NLG with 530B, and Google's Switch Transformer with 1.56 trillion parameters, to capture intricate patterns and relationships within the data, enabling them to generalize effectively across a wide array of tasks. 
Compared to these models, models used for robotic learning directly are relatively smaller, \eg RT-1~\cite{rt-1}, MOO~\cite{stone2023open}, and CLIPort~\cite{shridhar2022cliport} (with CLIP-ViT-Base), have 35M+, 111M, and 400M+, respectively. RT-2~\cite{rt-2} and variations of RT-2-X~\cite{9} incorporate pre-trained VLMs, which have large parameter space (\eg RT-2-PaLM-E and RT-2-PaLI-X have 12B and 55B, respectively).
Scaling up these models faces many challenges: 
\cib{1} maintaining real-time inference under strict time constraints (\eg 3-5Hz), 
\cib{2} maintaining robustness and coherence when generalizing to novel tasks and environments, 
\cib{3} addressing the computational demands and resource requirements of training and deployment, 
\cib{4} managing the complexity of integrating multimodal inputs, and 
\cib{5} ensuring reliability and safety in dynamic and unpredictable real-world settings.
As end-to-end techniques and integration of large vision, LLM, and VLM backbones become more prevalent~\cite{9} (see the evolution of robotic models in \autoref{fig:vla-timeline}), security research for these models is limited. Generalist robotic models continue to emerge, \eg OpenVLA-OFT, {GR00T}~N1, Gemini Robotics, pi-0.5, and WorldVLA~\cite{qu2025spatialvla,openvla_oft,gr00tn1,geminirobotics2025,pi05_2025,cen2025worldvla,cheang2024gr2}. Ensuring these models resilience to attacks is crucial.
To this end, we identify the following challenges.
% We identify the key challenges in ensuring model resilience to adversarial attacks.


\BfPara{C1: Wide-array of Data and Model Adoption}
The emergence of large-scale datasets for {\xo}-embodied generalist robots across diverse robotic platforms, embodiments, and environments has enabled many advances in the field.
Standardizing and processing these heterogeneous datasets in a consistent format requires substantial effort and innovation, as seen in efforts such as RoboNet~\cite{dasari2020robonet} and Open {\xo}-Embodiment Repository~\cite{9}. 
These datasets cover different kinds of tasks, objects, scenarios, and settings.
Moreover, recent work showed the benefits of learning visual representations from human video data~\cite{nair2022r3m,xiao2022masked,radosavovic2022,majumdar2024we,karamcheti2023language,mu2023ec2,bahl2023affordances} for robotics learning. 
Beyond robot-specific or {\xo}-embodied data, more data will inevitably be included to train generalist models.
This introduces security/safety challenges that must be addressed to prevent/mitigate vulnerabilities.


\BfPara{C2: Multimodal Data and Model Integration} 
The integration of multimodal data, \ie vision, language, and action, into large foundation models for robotics presents substantial security complexities. 
Despite the importance of this integration to improve robot functionality and autonomy through understanding and interacting with unseen environments, the heterogeneous nature of multimodal data greatly increases the attack surface, offering multiple avenues for malicious interference. 
Each data modality introduces particular vulnerabilities; for instance, adversarial images can fool vision models, while subtly changed text inputs can fool language models. Moreover, it is more challenging to determine the impact of adversarial manipulations on input data due to the varied ways in which such data and models are used in end-to-end robot learning.
We must identify and counteract attacks that can exploit the specific weaknesses of each modality. 
%Constructing effective defense mechanisms that can adjust to new threats and work across different modalities is no easy feat.

\BfPara{C3: Robustness and Trustworthiness}
Ensuring the robustness and trustworthiness of robotic agents is a major challenge. %, particularly given the methods used to train policies for robotic agents. 
Few studies have adequately assessed factors such as interpretability, robustness, safety, and fairness in robots, raising concerns about the potential for inherited vulnerabilities and bias from large backbone models. We must assess robot models
under a variety of conditions, including adversarial attacks and unexpected inputs, and develop methods that can improve their robustness.
We must develop methods for interpretability and fairness to ensure transparent and unbiased decisions.
%The complexity of training robotic policies, often involving reinforcement learning and continuous adaptation, complicates the assessment of these factors. 
%Without thorough evaluation, there is a risk that AI systems may inherit and propagate vulnerabilities from foundational models, compromising their safety and effectiveness. 

These challenges represent a tremendous opportunity for research into end-to-end learning models for {\xo}-embodied generalist robots. The PIs are well-positioned to advance this field through their expertise in AI security, robust model development, robotic systems, and distributed modeling and analysis.

\BfParaUL{Team}
\textbf{PI Abuhamad}
PI Abuhamad is an expert in AI security, especially in the security of large models. 
The PI has led significant work on secure and reliable AI systems~\cite{abuhamad2018large,abuhamad2020autosen,alasmary2020soteria,abuhamad2020multi,abusnaina2021dl,abuhamad2021large,abdukhamidov2023hardening,juraev2024unmasking,juraev2024impact,alawami2024motionid,abdukhamidov2024singleadv}, developing novel approaches to secure and interpretable AI \cite{abdukhamidov2023hardening,abusnaina2021dl,abdukhamidov2024singleadv,abdukhamidov2025stealthy,abuhamad2025nlp,abuhamad2025vit}, software security~\cite{abuhamad2018large,abuhamad2018large,abuhamad2020multi}, online security, privacy and safety~\cite{khan2023identification,reyes2023webtracker,shao2023lightweight,abuhamad2025shield}, and secure authentication methods~\cite{abuhamad2020autosen,alawami2024motionid}. 
These works appeared in ACM CCS, IEEE ICDCS, IEEE TDSC, ACM TOPS, IEEE TIFS, IEEE IoT-J, and PoPETS.  Because of this background and prior work, PI Abuhamad is uniquely positioned to carry out the proposed work.
\textbf{PI Thiruvathukal}
%TODO Add PI Thiruvathukal's relevant publications and contributions.

The PIs recognize the vast and evolving landscape of adversarial attacks against large foundation models. 
In this project, the PIs will work to establish foundational research and collaborative efforts to develop comprehensive frameworks and benchmarks that build toward secure and reliable robotic models. 
%While addressing all possible attack vectors is impossible within a single project, creating foundational tools, metrics, and methodologies will enable the broader research community to systematically evaluate and improve robotic system security. 
%The PI will leverage expertise in AI security to create lasting infrastructure that supports continuous improvement in robotic model security, with a particular focus on developing standardized evaluation protocols that can evolve alongside emerging threats.

